\documentclass[10pt,a4paper]{amsart}
\usepackage{amsmath,amssymb,amsthm,graphicx}
% Retained minimal preamble for the Stacks Project chapter sites-cohomology.tex (Cohomology on Sites).
% Verbatim excerpt of the chapter preamble: Xy-pic setup (native preamble lines 14-18) and the native
% theorem declarations with the mathematical macros (native preamble lines 154-217). Copied mechanically;
% nothing added, removed or redefined. The wrapper class and the amsmath/amssymb/amsthm/graphicx base
% packages are supplied by the assembler, so the class/externaldocument part of the native preamble is omitted.
% For commutative diagrams we use Xy-pic
\usepackage[all]{xy}

% We use 2cell for 2-commutative diagrams.
\xyoption{2cell}
% Theorem environments.
%
\theoremstyle{plain}
\newtheorem{theorem}[subsection]{Theorem}
\newtheorem{proposition}[subsection]{Proposition}
\newtheorem{lemma}[subsection]{Lemma}

\theoremstyle{definition}
\newtheorem{definition}[subsection]{Definition}
\newtheorem{example}[subsection]{Example}
\newtheorem{exercise}[subsection]{Exercise}
\newtheorem{situation}[subsection]{Situation}

\theoremstyle{remark}
\newtheorem{remark}[subsection]{Remark}
\newtheorem{remarks}[subsection]{Remarks}

\numberwithin{equation}{subsection}

% Macros
%
\def\lim{\mathop{\mathrm{lim}}\nolimits}
\def\colim{\mathop{\mathrm{colim}}\nolimits}
\def\Spec{\mathop{\mathrm{Spec}}}
\def\Hom{\mathop{\mathrm{Hom}}\nolimits}
\def\Ext{\mathop{\mathrm{Ext}}\nolimits}
\def\SheafHom{\mathop{\mathcal{H}\!\mathit{om}}\nolimits}
\def\SheafExt{\mathop{\mathcal{E}\!\mathit{xt}}\nolimits}
\def\Sch{\mathit{Sch}}
\def\Mor{\mathop{\mathrm{Mor}}\nolimits}
\def\Ob{\mathop{\mathrm{Ob}}\nolimits}
\def\Sh{\mathop{\mathit{Sh}}\nolimits}
\def\NL{\mathop{N\!L}\nolimits}
\def\CH{\mathop{\mathrm{CH}}\nolimits}
\def\proetale{{pro\text{-}\acute{e}tale}}
\def\etale{{\acute{e}tale}}
\def\QCoh{\mathit{QCoh}}
\def\Ker{\mathop{\mathrm{Ker}}}
\def\Im{\mathop{\mathrm{Im}}}
\def\Coker{\mathop{\mathrm{Coker}}}
\def\Coim{\mathop{\mathrm{Coim}}}

% Boxtimes
%
\DeclareMathSymbol{\boxtimes}{\mathbin}{AMSa}{"02}

%
% Macros for moduli stacks/spaces
%
\def\QCohstack{\mathcal{QC}\!\mathit{oh}}
\def\Cohstack{\mathcal{C}\!\mathit{oh}}
\def\Spacesstack{\mathcal{S}\!\mathit{paces}}
\def\Quotfunctor{\mathrm{Quot}}
\def\Hilbfunctor{\mathrm{Hilb}}
\def\Curvesstack{\mathcal{C}\!\mathit{urves}}
\def\Polarizedstack{\mathcal{P}\!\mathit{olarized}}
\def\Complexesstack{\mathcal{C}\!\mathit{omplexes}}
% \Pic is the operator that assigns to X its picard group, usage \Pic(X)
% \Picardstack_{X/B} denotes the Picard stack of X over B
% \Picardfunctor_{X/B} denotes the Picard functor of X over B
\def\Pic{\mathop{\mathrm{Pic}}\nolimits}
\def\Picardstack{\mathcal{P}\!\mathit{ic}}
\def\Picardfunctor{\mathrm{Pic}}
\def\Deformationcategory{\mathcal{D}\!\mathit{ef}}
\usepackage{cleveref}
\crefname{originalsection}{original section}{original sections}
\Crefname{originalsection}{Original section}{Original sections}
\makeatletter
\newlabel{derived-section-triangulated-sub}{{06UP}{source}{}{sourcepointer-1}{}}
\newlabel{derived-section-triangulated-sub@cref}{{[originalsection][0][]06UP}{[1][0][]source}{}{}{}}
\makeatother
\crefname{originallemma}{original lemma}{original lemmas}
\Crefname{originallemma}{Original lemma}{Original lemmas}
\makeatletter
\newlabel{sites-modules-lemma-chaotic-flat-quasi-coherent}{{0GZP}{source}{}{sourcepointer-2}{}}
\newlabel{sites-modules-lemma-chaotic-flat-quasi-coherent@cref}{{[originallemma][0][]0GZP}{[1][0][]source}{}{}{}}
\makeatother
\crefname{originallemma}{original lemma}{original lemmas}
\Crefname{originallemma}{Original lemma}{Original lemmas}
\makeatletter
\newlabel{sites-modules-lemma-chaotic-quasi-coherent}{{0GZN}{source}{}{sourcepointer-3}{}}
\newlabel{sites-modules-lemma-chaotic-quasi-coherent@cref}{{[originallemma][0][]0GZN}{[1][0][]source}{}{}{}}
\makeatother
\title{For a category with chaotic topology and flat restriction maps, quasi-coherent objects of the derived category are exactly the complexes with quasi-coherent cohomology}
\author{The Stacks Project Authors}
\date{stacks-project:sites-cohomology.tex:local-snapshot-20261011; GNU FDL 1.2 or later}
\begin{document}
\maketitle
\noindent\textit{Editorial note.} All delivered mathematical text is the author's own, copied from the pinned Stacks Project chapter file sites-cohomology.tex; no statement or proof step was rewritten, reordered or removed.\par\medskip
\noindent\textit{Editorial note.} The counted proof is the author's complete argument: the proof is written out in full: the definition of D\_QCoh(O) under the assumed flatness, the four-term computation of H\textasciicircum{}i(M)(U) for M quasi-coherent, and the converse obtained from the module-level criterion for quasi-coherent modules on a chaotic site.\par\medskip
\noindent\textit{Editorial note.} The label lemma-cartesian-flat-quasi-coherent names the counted statement, and the proof environment immediately following it in the source is delivered verbatim as its proof.\par\medskip
\noindent\textit{Editorial note.} The counted statement is carried by the author's own lemma environment in the chapter source.\par\medskip
\noindent\textit{Editorial note.} The statement and proof are copied from the pinned chapter file; the cited inputs are kept as external pointers with their canonical Stacks tags.\par\medskip
\noindent\textit{Editorial note.} The preamble is the chapter's own retained minimal preamble (sources/stacks-sites-cohomology/retained-preamble.tex), a verbatim excerpt of the native preamble: the Xy-pic setup and the native theorem declarations with the mathematical macros.\par\medskip
\noindent\textit{Allowed standard prerequisites (editorial).}\begin{itemize}
\item Chaotic topology and the identification of sections with values on objects. Same chapter, the chaotic topology on a category C and the exactness of taking sections over an object of C. Supplies the exactness of M |-> M(U) which yields the first and last equalities in the computation of H\textasciicircum{}i(M)(U).
\item Weak Serre property of the quasi-coherent modules and the module-level description. Modules on Sites, the lemmas stating that QCoh(O) is a weak Serre subcategory under the flatness assumption and characterising quasi-coherent modules on a chaotic site. Supplies the characterisation of quasi-coherence used in both directions of the proof.
\item Derived category of a weak Serre subcategory. Derived Categories, the section on triangulated subcategories, used to give meaning to D\_QCoh(O). Supplies the definition of the full subcategory D\_QCoh(O) appearing in the statement.
\end{itemize}\medskip
\ifdefined\hypertarget\hypertarget{sourcepointer-1}{}\fi
\noindent\textit{Source reference.} Original section 06UP: sources/stacks-chow/tags.txt, source lines 8753--8753. Prerequisite pointer invoked at proof line(s) 11572; it names a standard fact of a sibling chapter that the counted argument uses, and every step of the counted proof built on it is retained verbatim in the delivered body.\par\medskip
\ifdefined\hypertarget\hypertarget{sourcepointer-2}{}\fi
\noindent\textit{Source reference.} Original lemma 0GZP: sources/stacks-chow/tags.txt, source lines 21184--21184. Prerequisite pointer invoked at proof line(s) 11566; it names a standard fact of a sibling chapter that the counted argument uses, and every step of the counted proof built on it is retained verbatim in the delivered body.\par\medskip
\ifdefined\hypertarget\hypertarget{sourcepointer-3}{}\fi
\noindent\textit{Source reference.} Original lemma 0GZN: sources/stacks-chow/tags.txt, source lines 21183--21183. Prerequisite pointer invoked at proof line(s) 11591; it names a standard fact of a sibling chapter that the counted argument uses, and every step of the counted proof built on it is retained verbatim in the delivered body.\par\medskip
\begin{lemma}
\label{lemma-cartesian-flat-quasi-coherent}
Let $\mathcal{C}$ be a category viewed as a site with the chaotic
topology. Let $\mathcal{O}$ be a sheaf of rings on $\mathcal{C}$.
Assume for all $U \to V$ in $\mathcal{C}$ the restriction map
$\mathcal{O}(V) \to \mathcal{O}(U)$ is a flat ring map.
Then $\mathit{QC}(\mathcal{O})$ agrees with the subcategory
$D_\QCoh(\mathcal{O}) \subset D(\mathcal{O})$ of complexes
whose cohomology sheaves are quasi-coherent.
\end{lemma}

\begin{proof}
Recall that $\QCoh(\mathcal{O}) \subset \textit{Mod}(\mathcal{O})$
is a weak Serre subcategory under our assumptions, see
Modules on Sites, Lemma \ref{sites-modules-lemma-chaotic-flat-quasi-coherent}.
Thus taking the full subcategory
$$
D_\QCoh(\mathcal{O}) = D_{\QCoh(\mathcal{O})}(\textit{Mod}(\mathcal{O}))
$$
of $D(\mathcal{O})$ makes sense, see Derived Categories, Section
\ref{derived-section-triangulated-sub}. (Strictly speaking we don't need
this in the proof of the lemma.)

\medskip\noindent
Let $M$ be an object of $\mathit{QC}(\mathcal{O})$. Since for every
morphism $U \to V$ in $\mathcal{C}$ the restriction map
$\mathcal{O}(V) \to \mathcal{O}(U)$ is flat, we see that
\begin{align*}
H^i(M)(U)
& =
H^i(R\Gamma(U, M)) \\
& =
H^i(R\Gamma(V, M) \otimes_{\mathcal{O}(V)}^\mathbf{L} \mathcal{O}(U)) \\
& =
H^i(R\Gamma(V, M)) \otimes_{\mathcal{O}(V)} \mathcal{O}(U) \\
& =
H^i(M)(V) \otimes_{\mathcal{O}(V)} \mathcal{O}(U)
\end{align*}
and hence $H^i(M)$ is quasi-coherent by
Modules on Sites, Lemma \ref{sites-modules-lemma-chaotic-quasi-coherent}.
The first and last equality above follow from the fact that
taking sections over an object of $\mathcal{C}$ is an exact functor
due to the fact that the topology on $\mathcal{C}$ is chaotic.

\medskip\noindent
Conversely, if $M$ is an object of $D_\QCoh(\mathcal{O})$, then due to
Modules on Sites, Lemma \ref{sites-modules-lemma-chaotic-quasi-coherent}
we see that the map $R\Gamma(V, M) \to R\Gamma(U, M)$ induces isomorphisms
$H^i(M)(U) \to H^i(M)(V) \otimes_{\mathcal{O}(V)} \mathcal{O}(U)$.
Whence $R\Gamma(V, K) \otimes_{\mathcal{O}(V)}^\mathbf{L} \mathcal{O}(U)
\to R\Gamma(U, K)$ is an isomorphism in $D(\mathcal{O}(U))$ by the flatness
of $\mathcal{O}(V) \to \mathcal{O}(U)$ and we
conclude that $M$ is in $\mathit{QC}(\mathcal{O})$.
\end{proof}

\end{document}
